\documentclass[10pt,a4paper]{article}

\usepackage[a4paper,top=18mm,bottom=18mm,left=18mm,right=18mm]{geometry}
\usepackage{fontspec}
\usepackage{xcolor}
\usepackage{tabularx}
\usepackage{array}
\usepackage{enumitem}
\usepackage{tcolorbox}
\usepackage{fancyhdr}
\usepackage{lastpage}
\usepackage{titlesec}
\usepackage{microtype}
\usepackage{ragged2e}

\IfFontExistsTF{Arial}{
  \setmainfont{Arial}
  \setsansfont{Arial}
}{
  \setmainfont{TeX Gyre Heros}
  \setsansfont{TeX Gyre Heros}
}

\definecolor{AIEPRed}{HTML}{D62027}
\definecolor{AIEPNavy}{HTML}{102A43}
\definecolor{AIEPInk}{HTML}{1F2933}
\definecolor{AIEPMuted}{HTML}{536171}
\definecolor{AIEPLine}{HTML}{D7DEE8}
\definecolor{AIEPSoft}{HTML}{F6F8FB}
\definecolor{AIEPGold}{HTML}{B08A3C}
\definecolor{AIEPGreen}{HTML}{2F855A}

\pagestyle{fancy}
\fancyhf{}
\renewcommand{\headrulewidth}{0pt}
\fancyfoot[L]{\footnotesize\textcolor{AIEPMuted}{GeoGreen Escolar Osorno - Cronograma 2026}}
\fancyfoot[R]{\footnotesize\textcolor{AIEPMuted}{\thepage/\pageref{LastPage}}}

\setlength{\parindent}{0pt}
\setlength{\parskip}{5pt}
\hyphenpenalty=10000
\exhyphenpenalty=10000
\sloppy
\setlist[itemize]{leftmargin=*,topsep=2pt,itemsep=2pt}
\renewcommand{\arraystretch}{1.22}
\RaggedRight

\titleformat{\section}
  {\Large\bfseries\color{AIEPNavy}}
  {}{0pt}{}
\titleformat{\subsection}
  {\normalsize\bfseries\color{AIEPNavy}}
  {}{0pt}{}

\newcolumntype{Y}{>{\RaggedRight\arraybackslash}X}
\newcolumntype{L}[1]{>{\RaggedRight\arraybackslash}p{#1}}

\newtcolorbox{docbox}[2][]{
  colback=white,
  colframe=AIEPLine,
  boxrule=0.45pt,
  arc=1.2mm,
  left=3mm,
  right=3mm,
  top=2.4mm,
  bottom=2.4mm,
  before skip=5pt,
  after skip=5pt,
  title={\scriptsize\bfseries\MakeUppercase{#2}},
  coltitle=AIEPMuted,
  attach title to upper={\par\vspace{1mm}},
  #1
}

\newtcolorbox{accentbox}[2][]{
  colback=AIEPSoft,
  colframe=AIEPRed,
  boxrule=0.65pt,
  arc=1.2mm,
  left=3mm,
  right=3mm,
  top=2.4mm,
  bottom=2.4mm,
  before skip=6pt,
  after skip=6pt,
  title={\scriptsize\bfseries\MakeUppercase{#2}},
  coltitle=AIEPMuted,
  attach title to upper={\par\vspace{1mm}},
  #1
}

\newcommand{\eyebrow}[1]{%
  {\small\bfseries\textcolor{AIEPRed}{\MakeUppercase{#1}}}\par
  \vspace{8mm}
  {\color{AIEPRed}\rule{18mm}{0.8pt}}\par
  \vspace{10mm}
}

\newcommand{\datecell}[1]{\textbf{\textcolor{AIEPNavy}{#1}}}
\newcommand{\tag}[2]{%
  {\setlength{\fboxsep}{2pt}\colorbox{#1!10}{\scriptsize\bfseries\textcolor{#1}{\strut #2}}}%
}

\begin{document}

\eyebrow{Coordinación institucional}

{\Huge\bfseries\color{AIEPNavy}Cronograma post vacaciones\par}
\vspace{3mm}
{\Large\color{AIEPMuted}GeoGreen Escolar Osorno · Julio a octubre 2026\par}

\vspace{9mm}
\begin{tabular}{@{}>{\bfseries\color{AIEPNavy}}l@{\hspace{8mm}}l@{}}
Retorno de estudiantes & Lunes 27 de julio de 2026 \\
Hito ambiental & Viernes 2 de octubre de 2026 \\
Evento final sugerido & Lunes 5 de octubre de 2026 \\
Formato por sesión & Bloque A 09:00--10:30 · Bloque B 10:45--12:15 \\
\end{tabular}

\vspace{8mm}
\begin{accentbox}{Criterio de ajuste}
La semana del 27 de julio se reserva para reenganche y coordinación. La secuencia formativa comienza el lunes siguiente, lo que permite ampliar el período de mentorías y cerrar el proceso con un evento final asociado al hito ambiental de octubre.
\end{accentbox}

\begin{tabularx}{\textwidth}{@{}Y Y@{}}
\begin{docbox}[height=34mm]{Lógica pedagógica}
El programa avanza desde observar un problema ambiental local, comprender el material involucrado, imaginar una solución tecnológica y preparar una propuesta comunicable.
\end{docbox}
&
\begin{docbox}[height=34mm]{Lógica operativa}
Los talleres mantienen sus horarios originales. Las mentorías se distribuyen con semanas de avance y ajuste, evitando concentrar todo el trabajo en septiembre.
\end{docbox}
\end{tabularx}

\newpage
\section{Cronograma propuesto}

\begin{tabularx}{\textwidth}{@{}L{25mm}L{38mm}Y L{34mm}@{}}
\textbf{\color{AIEPNavy}Fecha} & \textbf{\color{AIEPNavy}Hito} & \textbf{\color{AIEPNavy}Propósito} & \textbf{\color{AIEPNavy}Producto} \\
\hline
\datecell{Lun 27 jul} & Reenganche y coordinación & Confirmar cursos, asistencia, sala, equipos y materiales. & Logística validada. \\
\datecell{Lun 03 ago} & Taller 1 & Observar residuos, hábitos y problemas del entorno. & Mapa y frase de problema. \\
\datecell{Lun 10 ago} & Taller 2 & Comprender materiales, reciclabilidad y separación en origen. & Ficha de material. \\
\datecell{Lun 17 ago} & Taller 3 & Presentar sensores, medición, alertas y GeoGreen como caso. & Idea tecnológica inicial. \\
\datecell{Lun 24 ago} & Mentoría 1 & Revisar problema, usuario y contexto afectado. & Problema validado. \\
\datecell{Lun 31 ago} & Semana de avance & Completar fichas, ajustar ideas y reunir evidencias. & Avance de equipo. \\
\datecell{Lun 07 sep} & Mentoría 2 & Definir solución, materiales, componentes o recursos. & Recursos definidos. \\
\datecell{Lun 14 sep} & Semana sin actividades & Respetar la semana de Fiestas Patrias. & Sin sesión programada. \\
\datecell{Lun 21 sep} & Mentoría 3 & Revisar maqueta, simulación, diagrama o prototipo. & Evidencia de avance. \\
\datecell{Lun 28 sep} & Mentoría 4 y ensayo & Afinar guion, soporte visual, roles y tiempos de pitch. & Pitch listo. \\
\datecell{Vie 02 oct} & Hito ambiental & Activar comunicación interna y convocatoria. & Difusión breve. \\
\datecell{Lun 05 oct} & Evento final & Presentar propuestas, recibir retroalimentación y reconocer equipos. & Presentaciones. \\
\datecell{Semana 12 oct} & Cierre interno & Consolidar evidencias, aprendizajes e informe final. & Informe y materiales. \\
\end{tabularx}

\newpage
\section{Lectura por etapa}

\begin{tabularx}{\textwidth}{@{}L{32mm}Y L{38mm}@{}}
\textbf{\color{AIEPNavy}Etapa} & \textbf{\color{AIEPNavy}Sentido} & \textbf{\color{AIEPNavy}Evidencia mínima} \\
\hline
Reenganche & Evita que el retorno post vacaciones absorba el primer taller y confirma condiciones reales. & Nómina, sala, horarios, materiales. \\
Talleres & Construyen la base común: problema, material y posibilidad tecnológica. & Fichas por equipo. \\
Mentorías & Transforman ideas iniciales en propuestas claras, viables y presentables. & Registro de avance por equipo. \\
Semanas de avance & Dan tiempo para completar trabajo, ordenar evidencias y preparar la siguiente sesión. & Ajustes y evidencia parcial. \\
Evento final & Cierra el proceso con presentación pública, retroalimentación y reconocimiento. & Pitch, fotografías, pauta o rúbrica. \\
Cierre interno & Ordena documentación para informe, continuidad y rendición del proyecto. & Informe final y set de materiales. \\
\end{tabularx}

\vspace{5mm}
\begin{accentbox}{Relación con octubre}
El viernes 2 de octubre funciona como hito comunicacional. Al caer viernes, el evento presencial se realiza el lunes 5 de octubre, manteniendo cercanía simbólica con la fecha y mejores condiciones operativas.
\end{accentbox}

\section{Criterios operativos}

\begin{tabularx}{\textwidth}{@{}Y Y@{}}
\begin{docbox}[height=40mm]{Mantener}
\begin{itemize}
  \item Los horarios por bloque.
  \item La secuencia de tres talleres.
  \item El trabajo por equipos.
  \item Evidencias simples y acumulativas.
\end{itemize}
\end{docbox}
&
\begin{docbox}[height=40mm]{Ajustar}
\begin{itemize}
  \item Fechas desde agosto.
  \item Mentorías con más espacio.
  \item Evento final hacia octubre.
  \item Semana posterior para cierre interno.
\end{itemize}
\end{docbox}
\end{tabularx}

\section{Resumen ejecutivo}

La propuesta desplaza el inicio formativo al lunes 3 de agosto, conserva los horarios acordados y extiende el proceso hasta octubre. El cambio evita usar el primer día post vacaciones como sesión crítica, mejora la distribución de mentorías y permite que el evento final dialogue con el hito ambiental del viernes 2 de octubre.

\end{document}
